\documentclass[11pt]{amsart}
\usepackage[margin=1in]{geometry}
\usepackage{amsmath,amssymb,amsthm,mathtools}
\usepackage[T1]{fontenc}
\usepackage{lmodern}
\usepackage{microtype}
\usepackage{enumitem}
\usepackage[hidelinks]{hyperref}

\newtheorem{theorem}{Theorem}[section]
\newtheorem{lemma}[theorem]{Lemma}
\newtheorem{proposition}[theorem]{Proposition}
\newtheorem{corollary}[theorem]{Corollary}
\theoremstyle{definition}
\newtheorem{definition}[theorem]{Definition}
\newtheorem{example}[theorem]{Example}
\theoremstyle{remark}
\newtheorem{remark}[theorem]{Remark}

\newcommand{\R}{\mathbb{R}}
\newcommand{\Z}{\mathbb{Z}}

\title[Window Neutrality and Phantom Balance]{Window Neutrality, Future-Balance Projection,\\
and the Phantom Balance Constraint}

\author{Jonathan Washburn}
\address{Austin, Texas, USA}
\email{washburn.jonathan@gmail.com}
\date{February 2026}

\begin{document}
\begin{abstract}
We study integer-valued signals on $\Z$ subject to a
\emph{window neutrality constraint}: the sum over every aligned
window of width~$W$ must vanish.  (We fix $W=8$ throughout, though the
results hold for any $W\ge2$.)  We prove three groups of results.
First, any nonzero contribution at a single tick forces a compensating
balance over the remaining ticks in its window
(the \emph{future-balance theorem}), with an explicit backward projection
bound $|d|\le B\cdot R$ relating the accumulated debt~$d$, the remaining
ticks~$R$, and the per-tick bound~$B$.
Second, we formalize the \emph{boundary-balance rule} (the last tick of
each window absorbs the residual) and prove it implies full window
neutrality by a telescoping argument.
Third, we prove that every 8-tick neutral schedule admits a unique
phase-correcting rotation offset that re-aligns it to any desired
boundary phase.
These results provide the windowed-cancellation and phantom-debt tools
used in companion papers on arithmetic phase bounds~\cite{F4} and
gauge-field defect annihilation.
\end{abstract}

\subjclass[2020]{Primary 39A10; Secondary 11B83, 68Q85}
\keywords{Discrete signal, window sum, neutrality constraint,
balance debt, backward projection, rotational invariance}
\maketitle

%=======================================================================
\section{Setup and definitions}\label{sec:setup}
%=======================================================================

Throughout, we fix the window width $W:=8$.  (Every result below holds
with $W$ replaced by any integer $\ge2$; only the specific numerics change.)

\begin{definition}[Window signal]\label{def:window}
A \emph{window signal} is a function $w:\{0,1,\ldots,W{-}1\}\to\Z$.
We say $w$ is \emph{neutral} if $\sum_{k=0}^{W-1} w(k)=0$.
\end{definition}

\begin{definition}[Lock event]\label{def:lock}
A \emph{lock event} at position $p\in\{0,\ldots,W{-}1\}$ with
contribution $\delta\in\Z\setminus\{0\}$ is an assignment
$w(p)=\delta$ for one tick of the window.
\end{definition}

\begin{definition}[Balance debt and remaining ticks]\label{def:debt}
Given a partial signal with contributions at positions $0,\ldots,k$
($0\le k\le W{-}1$), the \emph{balance debt} is $d:=\sum_{j=0}^k w(j)$
and the \emph{remaining ticks} is $R:=W{-}1{-}k$.
\end{definition}

%=======================================================================
\section{The future-balance theorem}\label{sec:future}
%=======================================================================

The first result says that in a neutral window, any single tick's
contribution determines an equal-and-opposite obligation on the rest.

\begin{theorem}[Lock forces future balance]\label{thm:lock}
Let $w$ be a neutral window signal.
If $w(p)=\delta$ for some $p\in\{0,\ldots,W{-}1\}$ with $\delta\ne 0$,
then
\[
\sum_{k\ne p} w(k)=-\delta.
\]
\end{theorem}
\begin{proof}
By neutrality, $\sum_{k=0}^{W-1} w(k)=0$.
Split the sum into $w(p)$ and the rest:
$w(p)+\sum_{k\ne p}w(k)=0$.
Hence $\sum_{k\ne p}w(k)=-w(p)=-\delta$.
\end{proof}

\begin{corollary}[Nonzero contribution creates debt]\label{cor:debt}
After a lock at position~$p$ with value $\delta\neq 0$, the remaining
$R=W{-}1{-}p$ ticks must collectively contribute $-\delta$ to achieve
window neutrality.
\end{corollary}
\begin{proof}
By Theorem~\ref{thm:lock}, the total over all other positions is $-\delta$.
Positions $0,\ldots,p{-}1$ are already fixed (their values are determined).
Hence the unfilled positions $p{+}1,\ldots,W{-}1$ (there are
$R=W{-}1{-}p$ of them) must supply the remaining $-\delta$ minus
whatever the earlier positions contributed.
In the worst case (earlier sum zero), the full obligation $-\delta$
falls on the remaining $R$ ticks.
\end{proof}

\begin{example}
With $W=8$ and a lock $w(2)=5$, positions $3,4,5,6,7$ (five remaining
ticks) must collectively contribute $-5$.
\end{example}

%=======================================================================
\section{Phantom magnitude and backward projection}\label{sec:phantom}
%=======================================================================

We now quantify how ``urgent'' the outstanding debt is.

\begin{definition}[Phantom magnitude]\label{def:phantom}
For balance debt $d\in\Z$ and remaining ticks $R\ge 0$:
\[
\Phi(d,R):=\frac{|d|}{R+1}.
\]
(The denominator $R+1$ rather than $R$ avoids division by zero when
$R=0$ and captures the interpretation that $R+1$ is the total number of
future decisions, including the current one.)
\end{definition}

\begin{proposition}[Phantom properties]\label{prop:phantom}
\begin{enumerate}[label=\textup{(\alph*)},nosep]
\item $\Phi(d,R)\ge 0$ for all $d,R$.
\item $\Phi$ is monotone increasing in $|d|$ for fixed~$R$.
\item $\Phi$ is monotone decreasing in~$R$ for fixed~$d\neq 0$.
\item $\Phi(0,R)=0$: zero debt implies zero phantom.
\item $\Phi(d,0)=|d|$: when no ticks remain, the phantom equals
  the full outstanding debt.
\end{enumerate}
\end{proposition}
\begin{proof}
All five properties follow directly from the formula
$\Phi(d,R)=|d|/(R+1)$ and the fact that $R+1\ge1>0$.
\end{proof}

\begin{theorem}[Backward projection principle]\label{thm:backward}
Suppose each of the $R$ remaining positions is bounded in magnitude
by~$B\ge 0$ (i.e.\ $|w(k)|\le B$ for each remaining position~$k$).
Then any feasible completion satisfying window neutrality requires
\[
|d|\le B\cdot R.
\]
Equivalently: if $|d|>BR$, no feasible completion exists.
\end{theorem}
\begin{proof}
Neutrality requires $\sum_{\text{remaining}}w(k)=-d$.
By the triangle inequality:
\[
|d|
=\Bigl|\sum_{\text{remaining}} w(k)\Bigr|
\le\sum_{\text{remaining}}|w(k)|
\le B\cdot R.
\qedhere
\]
\end{proof}

\begin{corollary}[Urgency bound]\label{cor:urgency}
Under the hypotheses of Theorem~\ref{thm:backward},
$\Phi(d,R)\le BR/(R+1)<B$.
\end{corollary}
\begin{proof}
By Theorem~\ref{thm:backward}, $|d|\le BR$.
Hence $\Phi(d,R)=|d|/(R+1)\le BR/(R+1)$.
Since $R/(R+1)<1$ for all $R\ge0$, we get $\Phi(d,R)<B$.
\end{proof}

\begin{remark}[Interpretation]
The urgency bound says the average debt per remaining decision is
strictly less than the per-tick capacity.  This means there is always
``room'' to service the debt---the phantom never exceeds the local budget.
\end{remark}

%=======================================================================
\section{Schedule neutrality and phase-correcting rotation}\label{sec:schedule}
%=======================================================================

We now extend from a single window to an infinite schedule on $\Z$.

\begin{definition}[Aligned $W$-tick schedule]\label{def:schedule}
An integer signal $s:\Z\to\Z$ satisfies \emph{$W$-tick neutrality} if
$\sum_{k=0}^{W-1} s(t_0+k)=0$ for every $t_0\equiv 0\pmod{W}$.
\end{definition}

\begin{definition}[Window index process]\label{def:idx}
Fix an initial phase $\iota_0\in\{0,\ldots,W{-}1\}$.
Define $\mathrm{idx}:\Z_{\ge 0}\to\{0,\ldots,W{-}1\}$ by
\[
\mathrm{idx}(0):=\iota_0,\qquad
\mathrm{idx}(n+1):=(\mathrm{idx}(n)+1)\bmod W.
\]
This is a deterministic cyclic counter with period $W$.
\end{definition}

\begin{theorem}[Window index tracking]\label{thm:index}
For the process in Definition~\ref{def:idx},
$\mathrm{idx}(n)=(\iota_0+n)\bmod W$ for every $n\ge0$.
\end{theorem}
\begin{proof}
By induction on~$n$.
\textit{Base case}: $\mathrm{idx}(0)=\iota_0=(\iota_0+0)\bmod W$. \checkmark

\textit{Inductive step}: assume $\mathrm{idx}(n)=(\iota_0+n)\bmod W$.
Then by Definition~\ref{def:idx}:
\[
\mathrm{idx}(n+1)
=(\mathrm{idx}(n)+1)\bmod W
=((\iota_0+n)\bmod W+1)\bmod W
=(\iota_0+n+1)\bmod W.
\qedhere
\]
\end{proof}

\begin{definition}[Boundary-balance rule]\label{def:boundary-balance}
An integer signal $s:\Z\to\Z$ satisfies the \emph{boundary-balance rule}
if for every block index $j\ge 0$:
\[
s(Wj+W{-}1)=-\sum_{k=0}^{W-2} s(Wj+k).
\]
That is, the last tick of each aligned block absorbs the negative of the
partial sum of the preceding $W{-}1$ ticks.
\end{definition}

\begin{theorem}[Neutrality from boundary balance]\label{thm:every8th}
If $s$ satisfies the boundary-balance rule, then $s$ satisfies
$W$-tick neutrality: for every $j\ge 0$,
\[
\sum_{k=0}^{W-1} s(Wj+k)=0.
\]
\end{theorem}
\begin{proof}
Fix $j\ge 0$.  Split the block sum:
\[
\sum_{k=0}^{W-1} s(Wj+k)
=\underbrace{\sum_{k=0}^{W-2} s(Wj+k)}_{=:P}+s(Wj+W{-}1).
\]
By the boundary-balance rule, $s(Wj+W{-}1)=-P$.
Hence the total is $P+(-P)=0$.
\end{proof}

\begin{theorem}[Schedule neutrality rotation]\label{thm:rotation}
Let $s$ satisfy $W$-tick neutrality.  Then there exists a unique
offset $r\in\{0,\ldots,W{-}1\}$ such that for every $k\ge 0$:
\[
\sum_{j=0}^{W-1} s(r+Wk+j)=0
\qquad\text{and}\qquad
\mathrm{idx}(r+Wk)=0.
\]
In particular, after reindexing by the offset $r$, neutrality holds on every
aligned block.  The offset is $r=(W-\iota_0)\bmod W$.
\end{theorem}
\begin{proof}
Set $r:=(W-\iota_0)\bmod W$.  By Theorem~\ref{thm:index},
$\mathrm{idx}(r)=(\iota_0+r)\bmod W=(\iota_0+W-\iota_0)\bmod W=0$.
Since $\mathrm{idx}(r)=0$, the $W$ consecutive ticks
$r,r{+}1,\ldots,r{+}W{-}1$ form a complete aligned block starting at a
multiple of $W$ in the reindexed coordinate.
By $W$-tick neutrality, their sum is zero.
Repeating with $r+Wk$ for every $k\ge0$ gives the full statement.

Uniqueness: if $r'$ also satisfies $\mathrm{idx}(r')=0$, then
$r'\equiv r\pmod{W}$, so $r'=r$ in $\{0,\ldots,W{-}1\}$.
\end{proof}

%=======================================================================
\section{Phantom cost augmentation}\label{sec:bridge}
%=======================================================================

When the window neutrality constraint is composed with a cost functional
(such as the $J$-cost of~\cite{F1}), the phantom balance acts as an
additive penalty that lifts the effective cost.

\begin{definition}[Augmented cost]\label{def:augmented}
For a configuration $x>0$, phantom data $(d,R)$, and
penalty scale $\lambda\ge 0$:
\[
J_{\mathrm{PL}}(x;d,R):=J(x)+\lambda\cdot\Phi(d,R).
\]
\end{definition}

\begin{proposition}[Augmented cost dominates base cost]\label{prop:augmented}
$J_{\mathrm{PL}}(x;d,R)\ge J(x)$ for all $x,d,R,\lambda\ge0$.
Equality holds iff $d=0$ or $\lambda=0$.
\end{proposition}
\begin{proof}
$\Phi(d,R)=|d|/(R+1)\ge0$ and $\lambda\ge0$, so
$\lambda\Phi\ge0$.  If $d\neq0$ and $\lambda>0$, then $\lambda\Phi>0$
and the inequality is strict.
\end{proof}

\begin{remark}[Phase-winding obstruction]
In applications to arithmetic phase dynamics (see~\cite{F4}),
the phantom mechanism forbids sustained one-sided phase winding:
any phase debt accumulated at tick~$k$ must be balanced by tick~$k{+}W{-}1$.
The backward projection bound (Theorem~\ref{thm:backward}) quantifies
the maximum sustainable debt at each intermediate tick.
Together, the future-balance theorem and backward projection provide a
discrete analog of the no-drift property: the windowed neutrality
constraint prevents unbounded accumulation.
\end{remark}

%=======================================================================
\section*{Acknowledgments}
%=======================================================================
The author thanks the referees for careful reading and for suggesting
the generalization to arbitrary window width~$W$.

\begin{thebibliography}{99}

\bibitem{F1}
J.~Washburn,
\emph{The canonical reciprocal cost: exact multi-bond minimization,
frustration lower bounds, and simultaneous-vs-sequential descent},
preprint, 2026.

\bibitem{F4}
J.~Washburn,
\emph{Phase caps, Herglotz positivity, and the Cayley--Schur pinch},
preprint, 2026.

\end{thebibliography}

\end{document}
